\chapter{Discrete Theta Terms}
\label{ch:global-discrete-theta-terms}

\ConventionsEuclidean

Discrete theta terms are topological phases in the path integral depending
on global bundle data rather than on local curvature polynomials alone.  They
are part of the definition of a gauge theory with a specified global form of
the gauge group.  Their correct formulation requires cohomology classes,
quadratic refinements, and compatibility with line-operator spectra.

\section{Finite Gauge Fields}

For a finite abelian group \(A\), a background \(A\)-gauge field on a
triangulated \(D\)-manifold \(M\) is represented by a cocycle
\[
  a\in Z^1(M,A)
\]
modulo coboundaries.  A discrete theta term is a \(U(1)\)-valued topological
action
\[
  \exp(2\pi\ii\int_M \omega(a)),
\]
where \(\omega\) is a cocycle representative of a class in
\[
  H^D(BA,U(1)).
\]
Changing \(\omega\) by a coboundary changes the action by a boundary term.
Thus on a closed manifold the phase depends only on the cohomology class and
the gauge-field isomorphism class.

\section{One-Form Backgrounds}

For an \(SU(N)/\mathbb Z_N\) gauge theory, the obstruction to lifting a
bundle to an \(SU(N)\)-bundle is a two-form cocycle
\[
  B\in Z^2(M,\mathbb Z_N).
\]
In four dimensions a discrete theta term is built from a quadratic refinement
of the cup square of \(B\).  On manifolds where the Pontryagin square
\(\mathcal P(B)\in H^4(M,\mathbb Z_{2N})\) is defined, the phase has the
form
\[
  \exp\!\left(
  \frac{2\pi\ii p}{2N}\int_M \mathcal P(B)
  \right),
  \qquad p\in\mathbb Z_{2N},
\]
with further identifications depending on spin structure and the precise
global form.  The integer \(p\) labels the discrete theta angle.

\section{Line Operators and Mutual Locality}

The choice of discrete theta term changes which electric and magnetic line
operators are mutually local.  Let a line charge be
\[
  \gamma=(e,m)\in \mathbb Z_N\oplus\mathbb Z_N .
\]
A discrete theta parameter \(p\) shifts the electric charge carried by a
magnetic line by a Witten-effect-type relation
\[
  e\sim p\,m\quad \text{mod }N
\]
in the minimal mutually local lattice.  Thus changing \(p\) is not a
decorative addition to the action; it changes the spectrum of admissible
genuine line operators.

\section{Coupling to Background Fields}

Let \(B\) be a background two-form field for a one-form symmetry.  The
generating functional \(Z[B]\) may transform under a one-form gauge
transformation \(B\mapsto B+d\lambda\) by a local phase.  If the phase cannot
be removed by a local counterterm in \(D\) dimensions, it is an anomaly
represented by an invertible theory in \(D+1\) dimensions.  Discrete theta
terms in \(D\) dimensions are therefore intertwined with anomaly inflow
terms in one higher dimension.

\begin{proposition}[Coboundary invariance on closed manifolds]
\label{prop:discrete-theta-coboundary-invariance}
If two discrete topological Lagrangians differ by a cocycle coboundary, their
exponentiated actions agree on closed manifolds for every flat finite
background field.
\end{proposition}

\begin{proof}
Let \(\omega'=\omega+d\alpha\).  For a closed \(D\)-manifold,
\[
  \int_M(\omega'-\omega)=\int_M d\alpha=\int_{\partial M}\alpha=0
  \quad \text{mod }1 .
\]
Therefore the exponentials of the two actions are equal.  On a manifold with
boundary the same calculation leaves the boundary phase, which must be
cancelled by boundary data or interpreted as the boundary variation of an
anomaly theory.
\end{proof}

\section{Continuum Gauge-Theory Status}

In continuum Yang--Mills theory with nontrivial global form, discrete theta
terms are defined by summing over bundles with specified characteristic
classes and weighting each topological sector by the corresponding phase.
A lattice regulator can implement the same data through plaquette variables
and higher-form background cochains.  The comparison between continuum and
lattice definitions requires matching characteristic classes and line
operator spectra.

\begin{openproblem}[Discrete theta terms in interacting continuum QFT]
\label{op:discrete-theta-continuum-qft}
Give a nonperturbative construction of four-dimensional gauge theories with
all global forms and discrete theta terms in which bundles, one-form
backgrounds, line operators, anomalies, and continuum limits are treated in a
single framework.
\end{openproblem}
